\section*{Addendum: Curve-Level Null-Refit Synthesis and the SN-RRF Family Interpretation}

Stage 26 provides a curve-level adversarial control for the SN-RRF extension. The locked radial-response kernel beats destructive nulls across the eligible population, but neighboring monotone saturating kernels compete. The resulting interpretation is not fixed-exponent uniqueness. It is a family-level diagnostic interpretation.

SN-RRF should be described as a scale-normalized residual-response family in which the locked $n=3.258993$ kernel functions as one representative kernel. The audit supports a radial-response shape signal relative to destructive nulls, but it does not establish that the original exponent is unique, fundamental, or physical.

\begin{itemize}
\item Valid population count: 143.
\item Stage 26 guard pass count: 5/5.
\item Median $\Delta\mathrm{AIC}$(RRF -- best destructive null): -85.0215.
\item RRF beats best destructive null fraction: 0.916.
\item RRF strongly beats destructive null fraction: 0.839.
\item Robust-class RRF beats destructive null fraction: 0.935.
\item Median $\Delta\mathrm{AIC}$(RRF -- best neighbor monotone kernel): 4.082.
\item RRF beats best neighbor monotone kernel fraction: 0.126.
\end{itemize}

This addendum does not support dark-matter exclusion, Lambda-CDM replacement, MOND/RAR defeat, Bullet-Cluster explanation, proof of a TVC physical mechanism, a Hubble-tension solution, or a universal fixed-exponent law.
